% Generated draft: requires mathematical and source review.
\documentclass[11pt]{article}
\usepackage[margin=25mm]{geometry}
\usepackage{fontspec}
\setmainfont{Latin Modern Roman}
\usepackage{amsmath,amssymb,mathtools,unicode-math}
\setmathfont{Latin Modern Math}
\usepackage{enumitem,needspace,xurl,hyperref}
\setlength{\parindent}{0pt}
\setlength{\parskip}{5pt}
\setlength{\emergencystretch}{3em}
\newcommand{\N}{\mathbb{N}}
\newcommand{\Z}{\mathbb{Z}}
\newcommand{\Q}{\mathbb{Q}}
\newcommand{\R}{\mathbb{R}}
\newcommand{\C}{\mathbb{C}}
\newcommand{\F}{\mathbb{F}}
\newcommand{\A}{\mathbb{A}}
\newcommand{\PP}{\mathbb{P}}
\newcommand{\E}{\mathbb{E}}
\newcommand{\eps}{\varepsilon}
\newcommand{\dd}{\,\mathrm{d}}
\begin{document}
\section*{UnsolvedMath ID: 700954}
\textbf{Source problem code:} catalogue entry
\subsection{Definitions and mathematical statement}
Let \(p\) be a prime. A pro-\(p\) group is a topological group that is an inverse limit of finite \(p\)-groups. Finite generation means topological finite generation. For a finitely generated pro-\(p\) group \(G\), let
\[
\mathbb F_p[[G]]=\varprojlim_{U\trianglelefteq_{\mathrm{open}}G}\mathbb F_p[G/U]
\]
be its completed group algebra, and let \(I\) be the kernel of its augmentation map to \(\mathbb F_p\). Define
\[
a_n(G)=\dim_{\mathbb F_p}\bigl(I^n/I^{n+1}\bigr),\qquad n\geq 1.
\]
The growth considered here is the growth of this sequence.

The supplied source does not specify an exact asymptotic definition of exponential growth or of growth slower than exponential. For an explicit provisional formulation, put
\[
\lambda(G)=\limsup_{n\to\infty}\bigl(\max\{1,a_n(G)\}\bigr)^{1/n},
\]
and interpret exponential growth as \(\lambda(G)>1\), and growth slower than exponential as \(\lambda(G)=1\). The latter condition is equivalent to
\[
\forall\varepsilon>0\ \exists N\geq1\ \forall n\geq N:\quad a_n(G)\leq e^{\varepsilon n}.
\]
These conventions require confirmation against the intended source terminology; they are not definitions explicitly supplied by the record.

Write \(H\leq_{\mathrm c}G\) when \(H\) is a closed subgroup of \(G\), and let \(F^{\mathrm{pro}\text{-}p}_2\) denote the free pro-\(p\) group on two generators. Isomorphisms of pro-\(p\) groups are understood to be continuous with continuous inverse.

\begin{enumerate}
\item[(a)] For every finitely generated pro-\(p\) group \(G\), does
\[
\lambda(G)>1
\quad\Longrightarrow\quad
\exists H\leq_{\mathrm c}G\text{ with }H\cong F^{\mathrm{pro}\text{-}p}_2
\]
hold?

\item[(c)] Does there exist a finitely generated pro-\(p\) group \(G\) such that
\[
\operatorname{cd}_p(G)<\infty,\qquad
G\text{ is not }p\text{-adic analytic},\qquad
\lambda(G)=1?
\]
Here \(\operatorname{cd}_p(G)\) is the continuous \(p\)-cohomological dimension: the least integer \(d\geq0\), if one exists, such that \(H^q_{\mathrm{cont}}(G,M)=0\) for every \(q>d\) and every discrete \(p\)-primary torsion continuous \(G\)-module \(M\). A \(p\)-adic analytic group is a group admitting a compatible finite-dimensional \(p\)-adic analytic manifold structure with analytic multiplication and inversion. The finite-generation requirement in this formulation is inherited from the opening definition of growth; part (c) itself does not repeat it. The source does not clarify whether \(p\) is prescribed arbitrarily or may be chosen as part of the existence assertion.

\item[(b), archived.] The supplied record also preserves the question attributed to A. Lubotzky and A. Shalev: for every finitely generated pro-\(p\) group \(G\), does
\[
\bigl(\exists H\leq_{\mathrm c}G:\ H\text{ is finitely generated and }\lambda(H)>1\bigr)
\quad\Longrightarrow\quad\lambda(G)>1
\]
hold? The archive supplied with the record reports a negative answer; this is a source-reported resolution, not an independently verified conclusion of this draft.
\end{enumerate}

For completeness, the supplied archived explanation takes \(G\) to be the Nottingham group over \(\mathbb F_p\), the group under composition of formal series \(t+\sum_{j\geq2}b_jt^j\), with \(b_j\in\mathbb F_p\), equipped with its natural coefficient topology. It reports that
\[
c_n(G)=\log_p\lvert G:\omega_n(G)\rvert
\]
has linear growth, where \(\{\omega_n(G)\}\) is the Zassenhaus filtration, and attributes the resulting subexponential growth of \(G\) to Quillen's theorem. It further reports that \(G\) contains a free pro-\(p\) subgroup of rank two, of exponential growth, citing R. Camina and a letter of M. Ershov. These claims and their applicability are retained as supplied archival background and require source review.
\subsection{Short English statement}
For a finitely generated pro-\(p\) group, measure growth by the dimensions of successive powers of the augmentation ideal in its completed group algebra. Part (a) asks whether exponential growth forces a free pro-\(p\) subgroup of rank two. Part (c) asks for finite cohomological dimension and growth slower than exponential without \(p\)-adic analyticity. The supplied archive reports a negative answer to part (b), concerning inheritance of exponential growth from a finitely generated closed subgroup. Exact growth conventions and the quantifiers in part (c) need confirmation.
\subsection{Sources}
\begin{itemize}
\item Local UnsolvedMath registry, permanent ID \(700954\), exact source problem code \texttt{catalogue entry}, title \emph{Kourovka Notebook Problem 13.36}; supplied complete source record, upstream revision \texttt{d371ee917ac9647fbb06cd6fc529deec6f3dda82}.
\item \emph{The Kourovka Notebook (Unsolved Problems in Group Theory)}, problem 13.36, proposed by O. V. Mel'nikov in the 13th issue (1995); the supplied record describes an update in the 21st issue (2026), \texttt{arXiv:1401.0300v48}. The supplied collection description names E. I. Khukhro and V. D. Mazurov as editors.
\item Archived part (b) is attributed in the supplied statement to A. Lubotzky and A. Shalev. Its supplied negative-answer explanation cites R. Camina, \emph{J. Algebra} \textbf{196} (1997), 101--113, and M. Ershov, \emph{Letter of 19.10.2009}. No article title or letter text is supplied.
\item The archived explanation invokes Quillen's theorem without a bibliographic reference or an exact theorem statement.
\item The supplied registry labels the problem \texttt{partially\_solved}. Its status review is dated \texttt{2026-10-16}, records evidence as \texttt{published}, and says that part (b) is in the Archive of Solved Problems while the remaining parts are in the main section. This is supplied dated curation, not an independently verified literature assessment.
\end{itemize}
\end{document}
